\section{Separators with a shared cutoff}\label{sec:separators}

Decomposing a normalized schedule at selected slots is a central device
in Dolev--Warmuth \citep[Theorem~2]{DW1984} and in the proper-decomposition
construction of Nederlof, Swennenhuis, and Węgrzycki
\citep[Section~3]{NSW2025}. Here we prove the shared-cutoff form needed
for the interval construction: successive boundaries must separate the
same ideal of low-priority jobs. The global-list invariants in later
sections control these descriptions inside descendant intervals.

This section concerns one fixed frame and a domain consisting of a
contiguous interval of a feasible schedule. A set is an \emph{ideal}
in the domain if it contains every domain predecessor of each of its
members. Its complement in the domain is an upset.

\begin{definition}\label{def:separator}
Given an ideal \(S\) of low jobs in a domain, a slot \(z\) is a
\emph{separator at cutoff \(S\)} if
\begin{enumerate}[label=(\alph*)]
 \item every domain job before \(z\) that is outside \(S\) belongs to \(P_z\);
 \item every domain job after \(z\) that belongs to \(S\) belongs to \(D_z\).
\end{enumerate}
The other domain jobs are called high. A left boundary outside the domain
is also a separator when all low domain jobs descend from it; there are
then no domain jobs before that boundary.
\end{definition}

An actual separator has a useful description of its earlier jobs. For
a global set \(H\subseteq J\), put
\begin{equation}
 B_z(H)=P_z\cup\bigl(H\setminus\wD z\bigr).
 \label{eq:B-definition}
\end{equation}

\begin{lemma}\label{lem:separator-description}
If \(z\) is an actual separator for a domain \(U\) at cutoff \(S\),
and \(H\cap U=S\), then \(B_z(H)\cap U\) is exactly the set of
domain jobs strictly earlier than \(z\). Reversing order and time,
and exchanging low with high, preserves the separator property.
\end{lemma}
\begin{proof}
Every job in \(P_z\) is strictly earlier than \(z\). An earlier low
job is outside \(\wD z\). A low job at \(z\) belongs to \(\wD z\),
and a later low job belongs to \(D_z\) by separation. Finally, every
earlier high job belongs to \(P_z\). These observations give both
inclusions. Reversal exchanges the two conditions in
\Cref{def:separator}.
\end{proof}

\subsection{Lexicographic normalization}

A \emph{priority order} is a total order of the domain extending
precedence. Read a full schedule chronologically, with each triple sorted
by priority. Among the finitely many feasible full schedules of the
domain, choose one whose resulting sequence of priority indices is
lexicographically least. Call it \emph{normalized} for that priority.
This is an existence choice, not an operation of the decision algorithm.

\begin{lemma}\label{lem:lex-prefix}
In a normalized schedule, a feasible exchange that puts a better-priority
job into an earlier slot is impossible. More generally, fix a slot \(r\).
After any feasible reordering confined strictly after \(r\), a feasible
exchange of a job at \(r\) with a better-priority job at a later slot
is still impossible.
\end{lemma}
\begin{proof}
Replacing one member of a sorted triple by a better-priority job
lexicographically improves that triple. In the second assertion, the
schedule through \(r\) is unchanged by the intervening reordering.
The proposed exchange would therefore first change the normalized
sequence at \(r\), and improve it there. The resulting full schedule
would contradict minimality, regardless of what happens later.
\end{proof}

\begin{lemma}[Shared-cutoff walk]\label{lem:shared}
Fix a normalized full schedule of a domain and let \(S\) and \(S^+\)
be consecutive priority prefixes, differing by the promotion of one job
from high to low. Suppose \(s\) is a separator at cutoff \(S\),
either inside the domain or at its exterior left boundary.
Then there is a separator \(t\ge s\), found by advancing only to
domain slots on the right, that separates both \(S\) and \(S^+\).
If the priority order is specified on a larger job set and the promoted
job is outside the domain, one may take \(t=s\).
\end{lemma}
\begin{proof}
\textbf{One advance at a fixed cutoff.}
Keep the normalized schedule and its priority order fixed. Let \(L\)
be a prefix of that priority order, and suppose the current position
\(s\) satisfies the earlier-high condition for \(L\).
If a later low job does not descend from \(s\), choose the earliest
slot \(t>s\) containing such a job, and choose one of them, \(x\).
Every domain predecessor of \(x\) lies strictly before \(s\).
Indeed, that predecessor is low because \(L\) is an ideal.
A predecessor at \(s\) would force \(x\in D_s\). A predecessor
strictly between \(s\) and \(t\) would descend from \(s\), by the
earliest choice of \(t\), and again force \(x\in D_s\).

Let \(y\) be a high domain job at a slot \(r\) with \(s\le r<t\).
If \(y\) had no domain successor at or before \(t\), exchanging
\(x\) and \(y\) would be feasible: all predecessors of \(x\) are
before \(s\), and delaying \(y\) would violate no successor
constraint. Moving \(x\) earlier and \(y\) later preserves their
other precedence constraints. Since \(L\) is a priority prefix,
\(x\) has better priority than \(y\), contradicting
\Cref{lem:lex-prefix}. Thus \(y\) has a successor by \(t\).
The successor is high, since the high set is an upset. If it is
strictly before \(t\), repeat the argument. Strictly increasing
times give a chain from \(y\) to a high job at \(t\).

When \(s\) is internal, every earlier high job has a descendant
at \(s\) by the earlier-high condition. That descendant is high,
and the preceding tracing extends the chain to \(t\).
When \(s\) is the exterior left boundary, there are no earlier
domain jobs. Hence \(t\) satisfies the earlier-high condition.
We may advance to \(t\) and repeat until there is no later low
exception. There are finitely many slots, and every advance is
strict, so the final position also satisfies the later-low condition.
Throughout this walk the normalized schedule itself stays fixed;
only the selected separator position changes.

\smallskip\noindent\textbf{Adding one low job.}
Apply this procedure with \(L=S^+\), starting at the given
separator for \(S\). The earlier-high condition for \(S^+\)
holds initially because its high set is smaller. If no advance is
needed, the same position separates both cutoffs.

Otherwise the first exception is the single promoted job: every
later job already low at \(S\) descends from the starting position.
At the first target, the promoted job lies in that target slot,
so it introduces no earlier high job for the old cutoff \(S\).
Thus the earlier-high condition holds there for both cutoffs.

In every subsequent advance, the exception is old-low: the promoted
job stays at the first target, at or before the current position.
An old-low exception has better priority than every old-high job.
The same exchange and successor tracing therefore propagates the
earlier-high condition for \(S\) as well as for \(S^+\).
At the final position, the later-low condition for \(S^+\) also
implies the later-low condition for \(S\), since \(S\subseteq S^+\).
The final position separates both cutoffs.

If the promoted job is outside the domain, the two restricted
prefixes coincide, so the original position already separates both.
\end{proof}

\begin{corollary}\label{cor:any-ideal}
Every ideal of a nonempty set with a full schedule has an actual
separator in some full schedule of that set.
\end{corollary}
\begin{proof}
Prioritize the ideal before its complement, using a linear extension
inside each. Normalize. The first actual slot separates the empty
prefix. Promote the jobs of the ideal one by one and use
\Cref{lem:shared}. The final position is an actual separator for
the ideal.
\end{proof}

The shared cutoff, rather than merely the existence of separate
separators, will be used for a gap: its low jobs descend from its
left endpoint, while its high jobs precede its right endpoint.
Keeping the final position of each promotion is important. Intermediate
positions assist the proof but need not become marks.
